% ======================================================================
% 02-case.tex — the scientific setting, condensed from TEX/formulation.tex
% and TEX/optimization.tex. Full derivations: see those files and Appendix A.
% ======================================================================
\section{The scientific setting}
\label{sec:case}

The project has two halves: neural surrogates for a contaminant-transport
equation, and a stochastic optimization layer that decides pumping rates under
geological uncertainty. Both are standard; we summarize only what the gates
need.

\subsection{One-dimensional sandbox and neural surrogates}
\label{sec:case-1d}

The verification problem is the dimensionless advection--dispersion equation
(ADE) for a conservative solute in a homogeneous, isotropic aquifer,
\begin{equation}
\pdd{u}{\tau}+\pdd{u}{\xi}
=\frac{1}{\mathrm{Pe}}\,\pddt{u}{\xi},
\qquad (\xi,\tau)\in(0,1)\times(0,\tau_{\max}],
\label{eq:ade}
\end{equation}
with step inlet and clean initial state,
\begin{equation}
u(\xi,0)=0,\qquad u(0,\tau)=1,\qquad u(1,\tau)=0,
\label{eq:bc}
\end{equation}
where $\xi=x/L$, $\tau=tv/L$ is measured in pore volumes, $u=C/C_0$, and
$\mathrm{Pe}=vL/D$. The closed-form Ogata--Banks solution
\citep{ogata1961solution}
\begin{equation}
u_{\mathrm{ref}}(\xi,\tau)=\frac12\left[
\erfc\!\left(\frac{(\xi-\tau)\sqrt{\mathrm{Pe}}}{2\sqrt{\tau}}\right)
+e^{\mathrm{Pe}\,\xi}\,
\erfc\!\left(\frac{(\xi+\tau)\sqrt{\mathrm{Pe}}}{2\sqrt{\tau}}\right)
\right]
\label{eq:ogata}
\end{equation}
is the ground truth. The shipped configuration uses $\mathrm{Pe}=20$,
$\tau_{\max}=0.5$ and a $256\times256$ test grid.

\paragraph{PINN.} A multilayer perceptron $u_\theta(\xi,\tau)$ ($\tanh$ hidden,
sigmoid output) minimizes
\begin{equation}
\mathcal{L}(\theta)=\lambda_{\mathrm{PDE}}\Lpde
+\lambda_{\mathrm{IC}}\Lic+\lambda_{\mathrm{BC}}\Lbc,
\label{eq:loss}
\end{equation}
with the PDE residual
$\mathcal{R}_\theta=\pdd{u_\theta}{\tau}+\pdd{u_\theta}{\xi}
-\mathrm{Pe}^{-1}\pddt{u_\theta}{\xi}$
formed by automatic differentiation and each term a mean square over resampled
collocation points. The best configuration uses a \emph{hard} ansatz that
satisfies the initial and outlet conditions identically and the inlet up to a
short ramp,
\begin{equation}
u_\theta(\xi,\tau)=(1-\xi)\,s(\tau)
+\xi(1-\xi)\,s(\tau)\,\mathrm{NN}_\theta(\xi,\tau),
\qquad s(\tau)=1-e^{-\tau/\varepsilon},
\label{eq:hardbc}
\end{equation}
with $\varepsilon=10^{-3}$. Both terms carry $s(\tau)$, so $u_\theta(\xi,0)=0$;
both carry $(1-\xi)$ or $\xi(1-\xi)$, so $u_\theta(1,\tau)=0$. Only the inlet is
approximate: $u_\theta(0,\tau)=s(\tau)\to1$ for $\tau\gg\varepsilon$.

\paragraph{Two artefacts worth knowing.} First, the step-inlet problem is
discontinuous at the corner $(\xi,\tau)=(0,0)$: the initial condition says
$u=0$, the inlet says $u=1$, and the reference returns $1$. Every continuous
ansatz that honors the initial condition returns $0$, so the pointwise error at
the corner is exactly one and a thin band around it is contaminated. This is a
metric artefact, not a network bug, and it is why two metrics are always
reported: the full-grid relative $L^2$ error and a corner-excluded interior
error over $\{\xi>0.05,\ \tau>0.05\}$. Second, the ramp contributes
$s'(0)=1/\varepsilon$, so
\begin{equation}
\max_{\xi,\tau}|\mathcal{R}_\theta|\gtrsim 1/\varepsilon
\label{eq:resfloor}
\end{equation}
\emph{independently of the weights}. At $\varepsilon=10^{-3}$ the corner
residual is $1.0\times10^{3}$ at initialization and after training, and on every
seed. A global residual maximum is therefore a useless diagnostic here; only
band means are meaningful.

\paragraph{FNO.} The Fourier neural operator \citep{li2020fourier} learns the
solution operator from parameters to the concentration field: a lifting MLP,
spectral convolution layers with $k_{\max}$ retained modes, and a projection
MLP. The 1D campaign uses a family of $(v,D)$ scenarios and evaluates against
\eqref{eq:ogata}. The architecture's input representation turns out to be the
whole story (\S\ref{sec:failures}).

\subsection{Two- and three-dimensional pump-and-treat}
\label{sec:case-opt}

The remediation problem couples steady Darcy flow to advective--dispersive
transport. With transmissivity $T=Kb$, head $h$, recharge $W$ and extraction
rates $q_i$ at wells $\xx_i$,
\begin{equation}
-\nabla\cdot\bigl(T(\xx)\nabla h\bigr)
=b\,W(\xx)-\sum_{i=1}^{n_w}q_i\,\delta_\varepsilon(\xx-\xx_i),
\label{eq:darcy}
\end{equation}
where $\delta_\varepsilon$ is a normalized Gaussian well kernel. Transport is
\begin{equation}
\phi R\,\pdd{C}{t}
=\nabla\cdot(\phi\mathbf{D}\nabla C)
-\nabla\cdot(\phi\mathbf{v}C)
-\sum_i q_i C(\xx_i)\delta_\varepsilon(\xx-\xx_i)+s(\xx,t),
\label{eq:transport}
\end{equation}
with velocity $\mathbf{v}=-K\nabla h/\phi$, dispersion tensor $\mathbf{D}$,
retardation $R$ and source $s$. Implicit Euler with first-order upwinding makes
the transport operator linear in $q$ and admits an exact discrete adjoint.

\paragraph{Random conductivity.} The primary random input is a log-normal field
\begin{equation}
K(\xx,\omega)=K_0\exp\!\bigl(G(\xx,\omega)\bigr),
\qquad
G=m_G+\sum_{n=1}^{N_{KL}}\sqrt{\lambda_n}\,\varphi_n(\xx)\,\xi_n(\omega),
\label{eq:kl}
\end{equation}
with $\xi_n\sim\mathcal{N}(0,1)$ and $(\lambda_n,\varphi_n)$ the eigenpairs of
an exponential covariance with variance $\sigma_G^2$ and correlation length
$\ell$ \citep{ghanem1991stochastic}. The benchmark uses
$\sigma_G=0.65$, $\ell=15$\,m, a 95\% variance truncation
($n_{KL}=1529$ at $160\times96$), nine candidate wells on a hexagonal lattice,
and a $10\ \mathrm{m^3/d}$ source at $100\ \mathrm{mg/L}$ for 3650 days.

\paragraph{The stochastic program.} The decision is a vector of nonnegative
pumping rates $\mathbf{q}$ in a convex admissible set $\mathcal{Q}$ (per-well
capacity, total budget, drawdown). The cost $J(\mathbf{q},\omega)$ combines
pumping energy, treatment, residual contamination and cleanup time. Compliance
is a random constraint $G(\mathbf{q},\omega)\le0$ at a monitoring line over the
final 20\% of the horizon, and the problem is
\begin{equation}
\min_{\mathbf{q}\in\mathcal{Q}}\ \E\bigl[J(\mathbf{q},\omega)\bigr]
\quad\text{s.t.}\quad
\Prob\bigl(G(\mathbf{q},\omega)\le0\bigr)\ge1-\alpha,
\label{eq:chance}
\end{equation}
with $\alpha=0.05$ and a head constraint handled analogously. The expectation
is approximated by sample-average approximation (SAA) over $M_{\text{train}}$
scenarios, and the chance constraint by a Rockafellar--Uryasev CVaR penalty
\citep{rockafellar2000optimization,shapiro2009lectures}. The search is a pattern
search with capacity projection and a batched compass sweep
\citep{audet2006mesh}. The full program, variants and algorithm discussion are
in \texttt{TEX/optimization.tex} and \texttt{TEX/SolvingAlgorithms.tex}.

\subsection{What counts as a result}
\label{sec:case-result}

The project uses four evidence classes, and every quantitative statement in
this article is labelled with one:
\begin{enumerate}
  \item \textbf{Derived}: true by construction (e.g.\ the $1/\varepsilon$
        residual floor), so seed-independent.
  \item \textbf{Multi-seed}: a median over $N\ge5$ seeds with its interquartile
        range, and a paired comparison for any variant.
  \item \textbf{Gated}: a run that clears a named threshold on a named artifact
        (mass balance, retained variance, backend parity, reliability).
  \item \textbf{Certified}: a policy re-evaluated out of sample, with a Wilson
        interval \citep{wilson1927probable}, not a bare fraction.
\end{enumerate}
The rest is \emph{preliminary} and is labelled as such. This taxonomy is the
article's unit of analysis: the mechanisms in \S\ref{sec:mechanisms} exist to
make each class cheap to produce, and the failures in \S\ref{sec:failures} are
cases where a weaker class was being reported as a stronger one.
